\documentclass[UTF8,a4paper,12pt]{ctexart}

\usepackage{xeCJK}
\usepackage{amsmath,amssymb,bm}
\usepackage{geometry}
\usepackage{booktabs}
\usepackage{array}
\usepackage{mathtools}
\usepackage{hyperref}
\usepackage{enumitem}

\geometry{left=2.6cm,right=2.6cm,top=2.5cm,bottom=2.5cm}
\hypersetup{
  colorlinks=true,
  linkcolor=blue,
  urlcolor=blue,
  citecolor=blue
}
\setlength{\parskip}{0.35em}
\setlength{\parindent}{2em}
\allowdisplaybreaks[2]
\renewcommand{\arraystretch}{1.2}

\title{\textbf{IR-SDE 方法介绍与数学推导}\\[0.6em]
\large 基于 CIFAR-10 条件图像恢复实现}
\author{周奕兵}
\date{}

\begin{document}

\begin{titlepage}
  \centering
  \vspace*{3cm}
  {\Huge\bfseries IR-SDE 方法介绍与数学推导\par}
  \vspace{1.2cm}
  {\Large 基于 CIFAR-10 条件图像恢复实现\par}
  \vspace{2.2cm}
  {\large 大作业第三部分\par}
  \vspace{0.6cm}
  {\Large 周奕兵\par}
  \vfill
  {\large
    本文从项目中的调度器、训练循环与 CIFAR-10 UNet 配置出发，\\
    推导前向随机过程、预测目标转换与反向跳步采样。\\[0.8em]
    编译方式：XeLaTeX
  \par}
  \vspace*{1.5cm}
\end{titlepage}

\tableofcontents
\clearpage

\section{基础变量定义与条件概率建模}

\subsection{图像、观测与扩散状态}

令一幅干净的 CIFAR-10 RGB 图像为
\[
  \bm{x}_0 \in [0,1]^{3\times H\times W},
\]
其中项目中的 CIFAR-10 图像分辨率为 \(H=W=32\)，图像经
\texttt{ToTensor()} 转换到 \([0,1]\) 数值范围。令
\[
  \bm{y}\in[0,1]^{3\times H\times W}
\]
表示与 \(\bm{x}_0\) 配对的含噪观测图像。数据加载器返回
\((\bm{y},\bm{x}_0)\)，训练时以 \(\bm{y}\) 为条件恢复 \(\bm{x}_0\)。

这里需要区分两种噪声。数据集构造 \(\bm{y}\) 时可以采用高斯噪声、
椒盐噪声、乘性散斑噪声或混合噪声，且最终会将观测裁剪到
\([0,1]\)。这些退化机制描述的是观测模型；它们不等同于下面
IR-SDE 前向过程中的高斯随机变量。IR-SDE 的扩散噪声由独立标准
高斯变量构造。

定义连续时间 \(t\geq 0\) 的扩散状态为 \(\bm{x}_t\)，其中
\(\bm{x}_{t=0}=\bm{x}_0\)。本实现只在整数时刻
\(t\in\{0,1,\ldots,T\}\) 取训练与推理状态。给定 \(\bm{x}_0\)
和观测 \(\bm{y}\)，其条件边缘分布为
\[
  q(\bm{x}_t\mid \bm{x}_0,\bm{y})
  =
  \mathcal{N}\!\left(
    \bm{y}+a_t(\bm{x}_0-\bm{y}),\,
    v_t\bm{I}
  \right).
\]
这里的协方差对所有像素通道使用同一标量方差 \(v_t\)，
\(\bm{I}\) 表示与图像张量维度相匹配的单位算子。

因此，模型的任务不是无条件地从纯噪声生成任意图像，而是学习
由观测 \(\bm{y}\) 条件化的恢复分布。训练时网络输入同时包含
当前扩散状态与观测；推理时整个时间序列也以固定的 \(\bm{y}\)
为条件。理想化地，可将恢复目标写作
\[
  p_\theta(\bm{x}_0\mid\bm{y}),
\]
而网络在给定时间状态下所学习的预测可表示为
\[
  f_\theta(\bm{x}_t,\bm{y},t).
\]
其中 \(\theta\) 在 \(p_\theta,f_\theta\) 中表示神经网络参数；
在 SDE 系数中，\(\theta>0\) 表示回复速率。两者语境不同。

\subsection{条件分布与像素空间直觉}

令
\[
  a_t=e^{-\theta t},\qquad
  v_t=\lambda^2\left(1-e^{-2\theta t}\right),\qquad
  \sigma_t=\sqrt{v_t}.
\]
则上面的条件分布等价于以下重参数化形式：
\begin{align}
  \bm{x}_t
  &=\bm{y}+a_t(\bm{x}_0-\bm{y})+\sigma_t\bm{\epsilon},
  \label{eq:reparam}
  \\
  \bm{\epsilon}&\sim\mathcal{N}(\bm{0},\bm{I}).
\end{align}
由此可见，前向过程的条件均值从 \(\bm{x}_0\) 出发，沿着
\(\bm{x}_0-\bm{y}\) 的方向逐渐缩小；随机项的方差则随时间增长。
因此，状态在平均意义上向观测 \(\bm{y}\) 靠拢，同时受到以
\(\bm{y}\) 为中心的高斯扰动。

\section{前向加噪过程：连续 SDE 与离散边缘分布}

\subsection{回复到观测图像的 Ornstein--Uhlenbeck 型 SDE}

考虑以 \(\bm{y}\) 为固定回复中心的随机微分方程
\begin{equation}
  d\bm{X}_t
  =
  \theta(\bm{y}-\bm{X}_t)\,dt
  +\sqrt{2\theta}\lambda\,d\bm{W}_t,
  \qquad
  \bm{X}_0=\bm{x}_0,
  \label{eq:sde}
\end{equation}
其中 \(\bm{W}_t\) 为与图像张量同维度的标准布朗运动，
\(\theta>0\) 控制向观测回复的速率，\(\lambda>0\) 控制长期噪声
标准差。令
\[
  \bm{Z}_t=\bm{X}_t-\bm{y}.
\]
因为 \(\bm{y}\) 对这条随机轨迹固定，\(d\bm{Z}_t=d\bm{X}_t\)，
故式 \eqref{eq:sde} 可改写为
\[
  d\bm{Z}_t=-\theta\bm{Z}_t\,dt
  +\sqrt{2\theta}\lambda\,d\bm{W}_t.
\]
两侧乘以积分因子 \(e^{\theta t}\)。由 It\^o 乘积公式，
确定性可微函数 \(e^{\theta t}\) 与随机过程相乘给出
\begin{align}
  d\!\left(e^{\theta t}\bm{Z}_t\right)
  &=\theta e^{\theta t}\bm{Z}_t\,dt
    +e^{\theta t}d\bm{Z}_t
  \\
  &=\theta e^{\theta t}\bm{Z}_t\,dt
    +e^{\theta t}
       \left(-\theta\bm{Z}_t\,dt
       +\sqrt{2\theta}\lambda\,d\bm{W}_t\right)
  \\
  &=\sqrt{2\theta}\lambda e^{\theta t}\,d\bm{W}_t.
\end{align}
从 \(0\) 积分到 \(t\)，并使用
\(\bm{Z}_0=\bm{x}_0-\bm{y}\)，得到
\begin{align}
  e^{\theta t}\bm{Z}_t
  &=\bm{x}_0-\bm{y}
    +\sqrt{2\theta}\lambda
      \int_0^t e^{\theta s}\,d\bm{W}_s,
  \\
  \bm{X}_t
  &=\bm{y}+e^{-\theta t}(\bm{x}_0-\bm{y})
    +\sqrt{2\theta}\lambda
      \int_0^t e^{-\theta(t-s)}\,d\bm{W}_s.
  \label{eq:ou-solution}
\end{align}
It\^o 随机积分是高斯随机变量，均值为零。利用 It\^o 等距公式，
其逐坐标方差为
\begin{align}
  \operatorname{Var}(\bm{X}_t\mid\bm{x}_0,\bm{y})
  &=2\theta\lambda^2
    \int_0^t e^{-2\theta(t-s)}\,ds\,\bm{I}
  \\
  &=2\theta\lambda^2
    \left[
      \frac{e^{-2\theta(t-s)}}{2\theta}
    \right]_{s=0}^{s=t}\bm{I}
  \\
  &=\lambda^2(1-e^{-2\theta t})\bm{I}.
\end{align}
同时，条件均值为
\[
  \mathbb{E}[\bm{X}_t\mid\bm{x}_0,\bm{y}]
  =\bm{y}+e^{-\theta t}(\bm{x}_0-\bm{y}).
\]
故将 \(a_t=e^{-\theta t}\) 和
\(v_t=\lambda^2(1-e^{-2\theta t})\) 代入，恰好得到
式 \eqref{eq:reparam} 的条件高斯边缘分布。

\subsection{精确离散转移与代码中的闭式加噪}

对任意时间间隔 \(\Delta>0\)，从式 \eqref{eq:ou-solution}
将起始时刻改为 \(t\)，可得精确转移：
\begin{align}
  \bm{X}_{t+\Delta}
  &=\bm{y}+e^{-\theta\Delta}(\bm{X}_t-\bm{y})
    +\sqrt{2\theta}\lambda
       \int_t^{t+\Delta}
         e^{-\theta(t+\Delta-s)}\,d\bm{W}_s
  \\
  &=\bm{y}+e^{-\theta\Delta}(\bm{X}_t-\bm{y})
    +\lambda\sqrt{1-e^{-2\theta\Delta}}\,\bm{\xi}_t,
  \label{eq:exact-transition}
\end{align}
其中 \(\bm{\xi}_t\sim\mathcal{N}(\bm{0},\bm{I})\)。第二行成立是因为
积分项为零均值高斯变量，其每个分量的方差为
\[
  2\theta\lambda^2\int_t^{t+\Delta}
      e^{-2\theta(t+\Delta-s)}\,ds
  =\lambda^2(1-e^{-2\theta\Delta}).
\]

当 \(\Delta=1\) 时，式 \eqref{eq:exact-transition} 给出相邻整数时刻
的离散转移。更重要的是，代码的 \texttt{add\_noise} 并不需要
逐步模拟这 \(t\) 次转移；它直接使用从初始状态到时刻 \(t\) 的
闭式边缘分布。代入 \(a_t,v_t\)，可逐项得到
\begin{align}
  \bm{x}_t
  &=\bm{y}+e^{-\theta t}(\bm{x}_0-\bm{y})
     +\sqrt{\lambda^2(1-e^{-2\theta t})}\,\bm{\epsilon}
  \\
  &=\bm{y}+a_t(\bm{x}_0-\bm{y})
     +\sqrt{v_t}\,\bm{\epsilon},
  \qquad \bm{\epsilon}\sim\mathcal{N}(\bm{0},\bm{I}).
  \label{eq:forward-code}
\end{align}
这就是实现中的加噪公式。它是线性变换加高斯随机变量的直接采样，
而不是 Euler--Maruyama 近似。代码使用
\(-\operatorname{expm1}(-2\theta t)\) 计算
\(1-e^{-2\theta t}\)，在 \(t\) 较小时比直接相减更有利于数值精度。

\subsection{边界时刻与有限时间的近似收敛}

在 \(t=0\) 时，
\[
  a_0=1,\qquad v_0=0,\qquad \bm{x}_{t=0}=\bm{x}_0.
\]
随着 \(t\) 增大，\(a_t\) 单调减小而 \(v_t\) 单调增大。若令
\(t\to\infty\)，则
\[
  a_t\to 0,\qquad v_t\to\lambda^2,
\]
因此
\[
  q(\bm{x}_t\mid\bm{x}_0,\bm{y})
  \longrightarrow
  \mathcal{N}(\bm{y},\lambda^2\bm{I}).
\]
但程序使用的是有限终点 \(T\)，所以严格地说
\[
  q(\bm{x}_T\mid\bm{x}_0,\bm{y})
  =
  \mathcal{N}\!\left(
    \bm{y}+a_T(\bm{x}_0-\bm{y}),v_T\bm{I}
  \right),
\]
只有当 \(a_T\) 足够小时，它才近似于
\(\mathcal{N}(\bm{y},\lambda^2\bm{I})\)。因此，推理初始化
\(\bm{x}_T=\bm{y}+\sqrt{v_T}\bm{\xi}\) 是以观测为中心的终点近似；
它并未显式加入条件均值中的 \(a_T(\bm{x}_0-\bm{y})\) 项。

项目默认令 \(T=100\)、\(\theta=5/T=0.05\)、\(\lambda=0.1\)。
此时 \(a_T=e^{-5}\)，而 \(v_T\) 已非常接近 \(\lambda^2\)，但
\(a_T\) 仍不是严格的零。系数随时间的变化见表
\ref{tab:schedule}。

\begin{table}[htbp]
\centering
\caption{默认参数下的前向过程系数（\(T=100,\theta=0.05,\lambda=0.1\)）}
\label{tab:schedule}
\begin{tabular}{ccccc}
\toprule
时间 \(t\) & \(a_t=e^{-\theta t}\) & \(v_t\) & \(\sigma_t=\sqrt{v_t}\)
& 条件均值 \\
\midrule
\(0\) & \(1.000000\) & \(0\) & \(0\)
& \(\bm{x}_0\) \\
\(50\) & \(0.082085\) & \(0.009933\) & \(0.099663\)
& \(\bm{y}+0.082085(\bm{x}_0-\bm{y})\) \\
\(100\) & \(0.006738\) & \(0.0099995\) & \(0.099998\)
& \(\bm{y}+0.006738(\bm{x}_0-\bm{y})\) \\
\bottomrule
\end{tabular}
\end{table}

\section{模型预测目标与变量转换}

\subsection{两种预测目标及其训练损失}

对于每个训练批次，程序对每个样本独立地从
\(\{1,\ldots,T\}\) 均匀采样一个整数时间步，再采样
\(\bm{\epsilon}\sim\mathcal{N}(\bm{0},\bm{I})\)，并依据
式 \eqref{eq:forward-code} 构造 \(\bm{x}_t\)。网络输入为
\((\bm{x}_t,\bm{y},t)\)，其输出记为
\[
  \bm{r}_\theta=f_\theta(\bm{x}_t,\bm{y},t).
\]
配置支持两类监督目标：

\begin{enumerate}[label=(\arabic*)]
  \item \textbf{sample 预测：}令网络直接预测干净图像，监督目标为
  \(\bm{x}_0\)；
  \item \textbf{epsilon 预测：}令网络预测前向过程中的标准高斯噪声，
  监督目标为 \(\bm{\epsilon}\)。
\end{enumerate}

两种模式使用的均方误差目标可统一写为
\begin{equation}
  \mathcal{L}(\theta)
  =
  \mathbb{E}_{(\bm{x}_0,\bm{y}),\,t,\,\bm{\epsilon}}
  \left[
    \left\|
      f_\theta(\bm{x}_t,\bm{y},t)-\bm{z}
    \right\|_2^2
  \right],
  \qquad
  \bm{z}=
  \begin{cases}
    \bm{x}_0, & \text{sample 模式},\\
    \bm{\epsilon}, & \text{epsilon 模式}.
  \end{cases}
  \label{eq:loss}
\end{equation}
程序默认使用 sample 模式，也允许通过配置选择 epsilon 模式。
均方误差训练下，理想的函数回归解是给定输入条件下目标变量的
条件期望；实际网络输出是这一回归目标的近似，并不意味着对单个
样本潜变量的精确反演。

\subsection{从 epsilon 预测反推干净图像}

从前向关系式
\[
  \bm{x}_t-\bm{y}
  =a_t(\bm{x}_0-\bm{y})+\sigma_t\bm{\epsilon}
\]
出发。若网络输出 \(\widehat{\bm{\epsilon}}\)，将其代入噪声位置，
并对 \(\bm{x}_0\) 做代数移项：
\begin{align}
  \bm{x}_t-\bm{y}
  &=a_t(\widehat{\bm{x}}_0-\bm{y})
    +\sigma_t\widehat{\bm{\epsilon}},
  \\
  a_t(\widehat{\bm{x}}_0-\bm{y})
  &=\bm{x}_t-\bm{y}
    -\sigma_t\widehat{\bm{\epsilon}},
  \\
  \widehat{\bm{x}}_0-\bm{y}
  &=\frac{
      \bm{x}_t-\bm{y}-\sigma_t\widehat{\bm{\epsilon}}
    }{a_t},
  \\
  \widehat{\bm{x}}_0
  &=\bm{y}
    +\frac{
      \bm{x}_t-\bm{y}-\sigma_t\widehat{\bm{\epsilon}}
    }{a_t}.
  \label{eq:epsilon-to-x0}
\end{align}
因为训练时间步满足 \(t\geq 1\)，且 \(a_t=e^{-\theta t}>0\)，
该除法在数学上有定义。随着 \(t\) 增大，\(a_t\) 变小，因此 epsilon
预测中的误差可能经除以 \(a_t\) 放大；这一代数事实也解释了为什么
不同预测目标的数值稳定性可能不同。

反过来，若已得到干净图像估计
\(\widehat{\bm{x}}_0\)，由同一前向式解出标准化噪声：
\begin{align}
  \bm{x}_t-\bm{y}
  &=a_t(\widehat{\bm{x}}_0-\bm{y})
    +\sigma_t\widehat{\bm{\epsilon}},
  \\
  \widehat{\bm{\epsilon}}
  &=\frac{
     \bm{x}_t-\bm{y}-a_t(\widehat{\bm{x}}_0-\bm{y})
    }{\sigma_t}.
  \label{eq:x0-to-epsilon}
\end{align}
当网络采用 sample 模式时，\(\bm{r}_\theta\) 即
\(\widehat{\bm{x}}_0\)，并可通过式 \eqref{eq:x0-to-epsilon}
求得配套的残差噪声。由于 \(\sigma_t>0\) 对所有 \(t\geq1\)
成立，该转换在代码所允许的正时间步上有定义。

上述过程主要是从高斯重参数化式出发的代数变换，而不是单独应用
贝叶斯公式求闭式后验。若从统计学习角度解释，MSE 训练会使网络
逼近给定输入下目标变量的条件均值；\texttt{prediction\_to\_x0\_noise}
则将该网络输出转换成调度器统一使用的
\((\widehat{\bm{x}}_0,\widehat{\bm{\epsilon}})\) 表示。

\begin{table}[htbp]
\centering
\caption{预测目标与统一变量转换}
\label{tab:prediction}
\begin{tabular}{p{0.17\linewidth}p{0.23\linewidth}p{0.48\linewidth}}
\toprule
模式 & 训练监督目标 & 推理时得到的变量 \\
\midrule
sample
& \(\bm{x}_0\)
& \(\widehat{\bm{x}}_0=\bm{r}_\theta\)，再用式
  \eqref{eq:x0-to-epsilon} 得到
  \(\widehat{\bm{\epsilon}}\)。 \\
epsilon
& \(\bm{\epsilon}\)
& \(\widehat{\bm{\epsilon}}=\bm{r}_\theta\)，再用式
  \eqref{eq:epsilon-to-x0} 得到
  \(\widehat{\bm{x}}_0\)。 \\
\bottomrule
\end{tabular}
\end{table}

\subsection{裁剪操作的含义与噪声重算}

调度器在执行采样步时可将干净图像估计裁剪至合法像素范围：
\[
  \widehat{\bm{x}}_0
  \leftarrow
  \Pi_{[0,1]}(\widehat{\bm{x}}_0)
  =
  \min\!\bigl(1,\max(0,\widehat{\bm{x}}_0)\bigr),
\]
其中 \(\Pi_{[0,1]}\) 表示逐像素投影到区间 \([0,1]\)。这利用了
训练图像的像素取值范围，防止网络对干净图像的估计落在无效范围。
它不是新的概率推断步骤，也不保证投影后的像素更接近真实值；
它是采样阶段的约束与数值后处理。

裁剪会改变 \(\widehat{\bm{x}}_0\)，所以原先网络输出的噪声一般不再
与新的 \(\widehat{\bm{x}}_0\) 满足同一个前向重构等式。代码因此在
裁剪之后重新计算
\[
  \widehat{\bm{\epsilon}}
  =
  \frac{
    \bm{x}_t-\bm{y}
    -a_t(\widehat{\bm{x}}_0-\bm{y})
  }{\sigma_t}.
\]
于是裁剪后的 \(\widehat{\bm{x}}_0\) 与重算的
\(\widehat{\bm{\epsilon}}\) 在当前 \(t\) 时刻仍严格满足
\[
  \bm{x}_t
  =\bm{y}
    +a_t(\widehat{\bm{x}}_0-\bm{y})
    +\sigma_t\widehat{\bm{\epsilon}}.
\]
在本项目中，训练阶段为诊断而调用变量转换时不启用该裁剪；
\texttt{step} 采样阶段则依据 \texttt{clip\_sample} 配置启用裁剪，
默认配置为开启。

\section{反向跳步采样机制}

\subsection{节点选择与初始化}

设推理节点为严格递减序列
\[
  T=t_0>t_1>\cdots>t_{K-1}\geq1,
\]
并在其后追加终点 \(t_K=0\)。代码可以按给定推理步数从
\(T\) 到 \(1\) 生成节点，也可以接受从 \(T\) 开始的自定义严格
递减整数节点。相邻推理循环变量分别是当前时刻
\(i=t_k\) 和前一推理时刻 \(j=t_{k+1}\)，因而允许 \(j<i-1\)，
即一次更新跨越多个训练时刻。

采样开始时，程序抽取
\(\bm{\xi}\sim\mathcal{N}(\bm{0},\bm{I})\)，并令
\begin{equation}
  \bm{x}_T=\bm{y}+\sqrt{v_T}\,\bm{\xi}.
  \label{eq:init}
\end{equation}
这对应于使用近似终点分布
\(\mathcal{N}(\bm{y},v_T\bm{I})\) 初始化。随后 UNet 每一步都接收
当前状态与固定观测的通道拼接
\(\operatorname{cat}[\bm{x}_i,\bm{y}]\)，并输出配置所指定的
sample 或 epsilon 预测。

\subsection{从当前时刻重构到较早时刻}

在当前时刻 \(i\)，调度器先把网络输出转换为一致的
\((\widehat{\bm{x}}_0,\widehat{\bm{\epsilon}})\)：
\begin{align}
  \bm{x}_i
  &=\bm{y}
    +a_i(\widehat{\bm{x}}_0-\bm{y})
    +\sigma_i\widehat{\bm{\epsilon}}.
  \label{eq:current-reconstruction}
\end{align}
在当前估计下，时间 \(j\) 的重构状态定义为
\begin{equation}
  \bm{x}_j
  =
  \bm{y}
  +a_j(\widehat{\bm{x}}_0-\bm{y})
  +\sigma_j\widehat{\bm{\epsilon}},
  \qquad j<i.
  \label{eq:reverse-step}
\end{equation}
式 \eqref{eq:reverse-step} 分别将同一干净图像估计与同一标准化噪声
按较早时刻的系数重新组合。它不是将当前 \(\bm{x}_i\) 原样缩放，
而是先分解当前状态中“相对观测的信号项”和“高斯噪声项”，再用
\(a_j,\sigma_j\) 重建。

当 \(j=0\) 时，
\[
  a_0=1,\qquad \sigma_0=0,
\]
故更新退化为
\[
  \bm{x}_0^{\mathrm{out}}
  =\bm{y}+(\widehat{\bm{x}}_0-\bm{y})
  =\widehat{\bm{x}}_0.
\]
也就是说，最后一步输出的是当前网络估计的干净图像；若采样时启用
裁剪，则这里的估计已经先被投影到 \([0,1]\)。

\subsection{跳步的性质与随机性来源}

式 \eqref{eq:reverse-step} 在一次更新中没有新采样随机变量，因此
给定初始噪声、观测图像与网络参数后，后续轨迹是确定的。随机性
来自式 \eqref{eq:init} 的初始噪声，而不是每个反向跳步额外注入的
噪声。对于稀疏推理节点，算法仍按所选 \(j\) 的边缘系数直接重构，
而不必逐个计算 \(i-1,i-2,\ldots,j\) 的中间状态。

这是一种基于预测干净样本与预测噪声的确定性跳步重构形式，结构上
与 DDIM 式更新相似。需要准确区分：它利用了前向过程的边缘关系
构造一致的较早时刻状态，但该操作本身并不等于对精确反向条件
分布 \(q(\bm{x}_j\mid\bm{x}_i,\bm{x}_0,\bm{y})\) 进行随机抽样。
实际效果还取决于网络对 \(\bm{x}_0\) 或 \(\bm{\epsilon}\) 的预测误差。

\subsection{有限值检查与输出裁剪}

CIFAR 推理循环在最终裁剪前先检查输出是否包含非有限值；若有
\(\mathrm{NaN}\) 或无穷值则报告错误，而不是将异常值静默转换为
看似有效的图像。检查通过后，pipeline 执行
\[
  \bm{x}_{\mathrm{out}}=
  \operatorname{clamp}(\bm{x},0,1).
\]
理由是数据集中的干净图像采用 \([0,1]\) 像素约定，而网络计算及
SDE 公式并不会自动保证最终预测落在该区间。最终裁剪使返回图像
满足合法像素范围，并适配后续图像显示与保存。应注意，最终输出
裁剪与调度器内部对 \(\widehat{\bm{x}}_0\) 的裁剪发生在不同位置：
前者约束 pipeline 的最终输出，后者会影响反向跳步中重算的噪声。

\begin{table}[htbp]
\centering
\caption{采样阶段各步骤及其数学作用}
\label{tab:sampling}
\begin{tabular}{p{0.18\linewidth}p{0.34\linewidth}p{0.38\linewidth}}
\toprule
阶段 & 计算 & 作用 \\
\midrule
初始化
& \(\bm{x}_T=\bm{y}+\sigma_T\bm{\xi}\)
& 以观测为中心构造近似终点状态。 \\
预测与转换
& \(f_\theta(\bm{x}_i,\bm{y},i)
  \mapsto(\widehat{\bm{x}}_0,\widehat{\bm{\epsilon}})\)
& 将两类训练目标统一成采样所需的变量对。 \\
跳步
& \(\bm{x}_j=\bm{y}+a_j(\widehat{\bm{x}}_0-\bm{y})
+\sigma_j\widehat{\bm{\epsilon}}\)
& 用较早时刻的系数重建状态，不在此步另抽噪声。 \\
结束
& \(j=0\)，随后输出 clamp
& 得到干净图像估计，并限制其像素范围。 \\
\bottomrule
\end{tabular}
\end{table}

\section{CIFAR UNet 结构与六通道输入}

\subsection{为何将 \(\bm{x}_t\) 与 \(\bm{y}\) 拼接}

CIFAR-10 图像有三个颜色通道。当前扩散状态
\(\bm{x}_t\in\mathbb{R}^{B\times3\times H\times W}\)，观测
\(\bm{y}\in\mathbb{R}^{B\times3\times H\times W}\)。项目沿通道维
拼接：
\begin{equation}
  \bm{u}_t
  =
  \operatorname{cat}[\bm{x}_t,\bm{y};\operatorname{dim}=1]
  \in\mathbb{R}^{B\times6\times H\times W}.
\end{equation}
因此输入为六通道，其中前三个通道是随时间变化的状态，后三个通道
是固定的条件观测。这一设计让卷积网络在每个空间位置都能同时访问
当前待恢复内容与对应的观测像素，并在多层卷积中学习两者的局部
对应关系。若只输入 \(\bm{x}_t\)，网络就没有显式获得该样本的
条件观测；通道拼接是将条件信息直接提供给 UNet 的一种简单实现。

这种拼接并不意味着模型把六个通道都当成待预测的图像。网络的
输出通道数为三，对应 RGB 图像形状；输出语义由调度器配置决定：
sample 模式下是三通道 \(\widehat{\bm{x}}_0\)，epsilon 模式下是
三通道 \(\widehat{\bm{\epsilon}}\)。

\subsection{实现中的 UNet 配置}

CIFAR UNet 的空间输入大小为 \(32\times32\)，
\texttt{in\_channels=6}、\texttt{out\_channels=3}，
\texttt{block\_out\_channels=(64,128)}，每个 block 配置
\texttt{layers\_per\_block=2}，归一化组数为 \(8\)。下采样路径和
上采样路径各配置两个 \texttt{DownBlock2D} 与
\texttt{UpBlock2D}。网络同时接收时间步 \(t\)，因此函数可写为
\[
  f_\theta:
  \mathbb{R}^{B\times6\times32\times32}\times t
  \longrightarrow
  \mathbb{R}^{B\times3\times32\times32}.
\]

从功能上看，下采样路径通过卷积特征变换和空间尺度压缩，逐步扩大
特征的有效感受野，使网络能够利用更大范围的图像上下文；上采样路径
逐步恢复空间分辨率，将高层语义信息映射回像素尺度。UNet 的跳跃
连接将较浅层的高分辨率特征传递到对应的上采样阶段，有助于同时
保留边缘、纹理等局部细节与深层语义信息。时间步嵌入使相同的卷积
网络能够根据当前噪声尺度改变预测行为。

\begin{table}[htbp]
\centering
\caption{CIFAR-10 条件 UNet 的关键结构配置}
\label{tab:unet}
\begin{tabular}{lll}
\toprule
配置项 & 实际取值 & 对应含义 \\
\midrule
输入分辨率 & \(32\times32\) & CIFAR-10 的空间尺寸 \\
输入通道 & \(6\) & \(3\) 通道 \(\bm{x}_t\) 加 \(3\) 通道 \(\bm{y}\) \\
输出通道 & \(3\) & RGB 的 sample 或 epsilon 预测 \\
Block 通道宽度 & \((64,128)\) & 两级特征通道配置 \\
每个 Block 层数 & \(2\) & 每个阶段的层数配置 \\
归一化组数 & \(8\) & Group Normalization 配置 \\
下采样/上采样类型
& \texttt{DownBlock2D}/\texttt{UpBlock2D}
& 编码与解码路径 \\
\bottomrule
\end{tabular}
\end{table}

\section{方法流程总结与实现对应关系}

将训练与推理合并来看，IR-SDE 的关键并非在每一时刻重新独立预测
一张无条件图像，而是围绕同一个观测 \(\bm{y}\) 建立由干净图像到
观测邻域的前向过程，再学习在不同噪声尺度下从状态中恢复干净
图像信息。训练时对每张配对样本执行
\[
  (\bm{x}_0,\bm{y})
  \longrightarrow
  t,\bm{\epsilon}
  \longrightarrow
  \bm{x}_t
  \longrightarrow
  f_\theta(\bm{x}_t,\bm{y},t)
  \longrightarrow
  \mathcal{L};
\]
推理时则执行
\[
  \bm{y}
  \longrightarrow
  \bm{x}_T
  \longrightarrow
  (\widehat{\bm{x}}_0,\widehat{\bm{\epsilon}})
  \longrightarrow
  \bm{x}_{t_{k+1}}
  \longrightarrow\cdots\longrightarrow
  \widehat{\bm{x}}_0.
\]

综上，代码中的核心数学对象可以归纳为以下四点：
\begin{enumerate}[label=\arabic*.]
  \item 前向 SDE 的回复中心是观测图像 \(\bm{y}\)，其解析边缘分布
  使任意训练时间步的 \(\bm{x}_t\) 都可一次闭式采样。
  \item 网络可学习预测 \(\bm{x}_0\) 或标准高斯噪声
  \(\bm{\epsilon}\)；调度器通过代数变换将两种输出统一为
  \((\widehat{\bm{x}}_0,\widehat{\bm{\epsilon}})\)。
  \item 反向跳步以较早时刻的 \(a_j,\sigma_j\) 重构状态，
  不需要按训练步数逐步遍历，也不在每个跳步重新注入随机噪声。
  \item CIFAR UNet 将当前状态与观测拼接成六通道条件输入，
  并输出三通道预测；最终结果再限制在 \([0,1]\) 像素区间。
\end{enumerate}

\end{document}
